\documentclass[12pt]{article}
\usepackage[utf8]{inputenc}
\usepackage{amsmath, amssymb}
\usepackage{geometry}
\geometry{margin=1in}

\title{Solving for $L_M$ Step-by-Step}
\author{Specific-Factors Model Example}
\date{}

\begin{document}
\maketitle

\section*{Setup}

We have the production functions:
\[
Q_M = K^{1/3}L_M^{2/3}, \qquad Q_F = T^{1/2}L_F^{1/2}, \qquad L_F = 90 - L_M.
\]

Given endowments:
\[
K=64 \;(\Rightarrow K^{1/3}=4), \qquad T=100 \;(\Rightarrow T^{1/2}=10).
\]

The equilibrium condition for labor mobility is:
\[
P_M \cdot MPL_M = P_F \cdot MPL_F.
\]

The marginal products are:
\[
MPL_M = \frac{2}{3}K^{1/3}L_M^{-1/3} = \frac{8}{3}L_M^{-1/3}, \qquad
MPL_F = \frac{1}{2}T^{1/2}L_F^{-1/2} = 5(90-L_M)^{-1/2}.
\]

---

\section*{Case 1: Initial Prices $P_M=2$, $P_F=1$}

\noindent Step 1: Substitute into the equilibrium condition
\[
2 \cdot \frac{8}{3}L_M^{-1/3} = 1 \cdot 5(90 - L_M)^{-1/2},
\]
which simplifies to
\[
\frac{16}{3}L_M^{-1/3} = 5(90 - L_M)^{-1/2}.
\]

\bigskip
\noindent Step 2: Take reciprocals to make exponents positive
\[
\frac{3}{16}L_M^{1/3} = \frac{1}{5}(90 - L_M)^{1/2}.
\]

\bigskip
\noindent Step 3: Clear constants and square both sides to eliminate the square root
\[
\left(\frac{15}{16}\right)^2 L_M^{2/3} = (90 - L_M),
\]
or equivalently,
\[
\frac{225}{256}L_M^{2/3} = 90 - L_M.
\]

\bigskip
\noindent Step 4: Substitute $y = L_M^{1/3}$ so that $L_M = y^3$ and $L_M^{2/3} = y^2$
\[
\frac{225}{256}y^2 = 90 - y^3 \quad \Rightarrow \quad y^3 + \frac{225}{256}y^2 - 90 = 0.
\]

\bigskip
\noindent Step 5: Solve the cubic numerically  
Testing values:

\[
\begin{aligned}
y=4.20 &\Rightarrow f(y) = 4.20^3 + \frac{225}{256}(4.20)^2 - 90 \approx -0.42,\\
y=4.22 &\Rightarrow f(y) \approx +0.81.
\end{aligned}
\]
By interpolation, $y \approx 4.2068$ gives $f(y) \approx 0$.

\bigskip
\noindent Step 6: Compute $L_M$
\[
L_M = y^3 \approx (4.2068)^3 \approx 74.45.
\]
Hence:
\[
\boxed{L_M = 74.45, \quad L_F = 90 - 74.45 = 15.55.}
\]

---

\section*{Case 2: After Price Increase $P_M=3$, $P_F=1$}

\noindent Step 1: New wage condition
\[
3 \cdot \frac{8}{3}L_M^{-1/3} = 5(90 - L_M)^{-1/2} \quad \Rightarrow \quad 8L_M^{-1/3} = 5(90 - L_M)^{-1/2}.
\]

\bigskip
\noindent Step 2: Take reciprocals and square
\[
\left(\frac{5}{8}\right)^2 L_M^{2/3} = 90 - L_M,
\]
\[
\frac{25}{64}L_M^{2/3} = 90 - L_M.
\]

\bigskip
\noindent Step 3: Substitute $y = L_M^{1/3}$ again
\[
y^3 + \frac{25}{64}y^2 - 90 = 0.
\]

\bigskip
\noindent Step 4: Solve numerically  
\[
\begin{aligned}
y=4.35 &\Rightarrow f(y) \approx -0.28,\\
y=4.36 &\Rightarrow f(y) \approx +0.33.
\end{aligned}
\]
The root is near $y = 4.3534$.

\bigskip
\noindent Step 5: Compute $L_M$
\[
L_M = y^3 \approx (4.3534)^3 \approx 82.59,
\]
so
\[
\boxed{L_M = 82.59, \quad L_F = 90 - 82.59 = 7.41.}
\]

---

\section*{Summary of Results}
\[
\begin{array}{lcc}
\text{Case} & L_M & L_F \\
\hline
P_M = 2 & 74.45 & 15.55 \\
P_M = 3 & 82.59 & 7.41 \\
\end{array}
\]

\noindent As $P_M$ rises, labor moves into manufacturing, confirming the positive link between the relative price of $M$ and $L_M$ predicted by the specific-factors model.

\end{document}
